%======= Einbinden der benötigten Packete (= Zusatzfunktionen)
\usepackage[T1]{fontenc}                % für Fonts in westeuropäischer Codierung
\usepackage{lmodern}										% Latin Modern Paket verändert die verwendete Schriftart. Bessere Darstellung für pdf
\usepackage{textcomp}										% Provides extra symbols, e.g. arrows like \textrightarrow, various currencies (\texteuro,..
\usepackage[utf8]{inputenc}    				% german special characters
\usepackage[english]{babel}     % hyphenation   usepackage[english,ngerman]{babel} f. deutsch
\usepackage{pdfpages}										% Einbinden von pdf-Files
\usepackage{pifont,textcomp,mathcomp}   % dingbad psfonts and text-compilant fonts (euro, TM, ...)
\usepackage{amsmath,amsopn,amsthm}      % AMS mathematics
\usepackage{amssymb}                    % zusätzliche Symbole 
\usepackage{xspace}											% avoids eaten spaces

%======= Eine Umgebung für Bilder und Tabellen			  	
%\usepackage[textfont={Small},labelfont={bf},margin=1cm,format=plain,font=singlespacing]{caption}   % hanging caption text [hang]
\usepackage[textfont={small},labelfont={bf},margin=1cm,format=plain,font=singlespacing]{caption}   % hanging caption text [hang]
\captionsetup*[figure]{name=Fig.}			 % Abbildungsunterschrift beginnt mit Abb.
\captionsetup*[table]{name=Tab.}


%======= Farben für Überschriften
\usepackage{color}
\definecolor{TUBlue}{rgb}{0,0.4,0.6}   % TU blue RGB 0 102 153
\addtokomafont{sectioning}{\sffamily\bfseries\selectfont\color{TUBlue}}
\setkomafont{chapter}{\normalfont\huge\sffamily\bfseries\color{TUBlue}}
\addtokomafont{section}{\Large}
\addtokomafont{subsection}{\normalfont\Large\sffamily\bfseries\color{TUBlue}}
\addtokomafont{subsubsection}{\normalfont\large\sffamily\bfseries\color{TUBlue}}
\addtokomafont{paragraph}{\normalfont\large\sffamily\bfseries\color{TUBlue}}


%======= Kopf-/Fußzeilen
\pagestyle{plain} % nur Fußzeile



%======= Eine kompakte Umgebung für die Bilder
\newcommand{\bild}[5]{{
\begin{figure}[#1]
\begin{center}
\includegraphics[width=#2\linewidth]{#3}
\caption{#4}\label{#5}
\end{center}
\end{figure}
}}

\newcommand{\bilder}[9]{{
    \begin{figure}[#1]
    \begin{center}
    \begin{subfigure}{0.47\textwidth}
    \includegraphics[width=#2\linewidth]{#3}
    \caption{#4}
    \end{subfigure}
    \begin{subfigure}{0.47\textwidth}
    \includegraphics[width=#5\linewidth]{#6}
    \caption{#7}
    \end{subfigure}
    \caption{#8}\label{#9}
    \end{center}
    \end{figure}
}}

\newcommand{\tabelle}[5]{{
\begin{table}[#1]
    \centering
    \caption{#2}
    \label{#3}
    \begin{tabular}{#4}
        #5
    \end{tabular}
\end{table}
}}

%======= Definitionen eigener Befehle
